% Research candidate only; does not edit the main manuscript.
% Coordinates refer to the two adjacent restored vertices. The available
% lists are three-element subsets of a labelled four-colour palette.

\begin{lemma}[Common edges for arbitrary adjacent-pair fibre cycles]
Let $F(A,B)$ consist of pairs $(x,y)$ with $x\in A$, $y\in B$, $x\ne y$,
with one-coordinate changes as edges, where $A,B$ are three-element
subsets of a four-element palette. If one available list changes and
the other is fixed, every Hamilton cycle of the old fibre and every
Hamilton cycle of the new fibre have at least two common edges.
The assertion also holds for an unchanged fibre.
\end{lemma}
\begin{proof}
For an unchanged fibre with equal lists the graph is a $6$-cycle.
With unequal lists its two states using the two common colours have
degree two, giving four distinct forced edges in every Hamilton cycle.

By exchanging coordinates, suppose only $A$ changes. If either its
old or new value equals $B$, write the common first list as $\{a,b\}$
and $B=\{a,b,c\}$. The common fibre is the path
\[
ab,\ ac,\ bc,\ ba.
\]
Its two end edges are forced in both Hamilton cycles, since $ab,ba$
have degree two in both full fibres.

Otherwise, write the old and new excluded colours as $a,b$, and
the fixed excluded colour as $c$, with fourth colour $d$. The common
fibre has vertices $ca,cb,cd,da,db$ and six edges. In the old fibre,
the two outside states are $ba,bd$. The state $bd$ has degree two,
forcing the consecutive path $ba,bd,cd$ in every Hamilton cycle.
Consequently the cycle has exactly four edges in the common fibre.
The new fibre has the analogous outside path $ab,ad,cd$ and also
retains four common-fibre edges. Their intersection therefore has
at least $4+4-6=2$ edges.
\end{proof}

% Potential use: the adjacent-two-four proof needs only one displayed
% 7-state unequal-list cycle, rather than choosing between two templates.
% This is a stronger local overlap property, not a higher-palette or
% original-graph threshold claim. An actual external move changes at most
% one of the lists because the two exterior neighbours are distinct.
